\chapter{Wound Suturing}

\section*{Overview}
Wound suturing involves stitching the edges of a laceration or surgical incision together to promote healing by primary intention, minimize scarring, and reduce infection risk.

\section*{Alternative Names}
Suture repair, Stitches

\section*{Principle}
The skin edges of a clean wound are approximated with sutures that hold the tissue layers in alignment during healing. Closure reduces wound tension, controls bleeding, and restores tissue continuity, promoting healing by primary intention while minimizing scarring and infection risk.
\section*{History}
Suturing was practiced in ancient Egypt using plant fibers and hair. Joseph Lister introduced sterile catgut sutures in the 1860s. Synthetic polymers were introduced in the 20th century.

\section*{Why It Is Done}
Ordered to close open lacerations or wounds that are deep, have gaping edges, or are bleeding persistently.

\section*{Is This a Primary or Secondary Test or Procedure}
Primary definitive treatment for skin lacerations requiring primary closure.

\section*{How It Is Performed}
The wound is irrigated and cleaned. Local anesthetic is injected. The clinician uses a needle holder, suture material, and tissue forceps to sew the wound edges together.

\section*{Preparation}
Irrigate the wound thoroughly with normal saline. Debride devitalized tissue if necessary.

\section*{Contraindications and Precautions}
Wounds that are heavily contaminated, showing signs of active infection, or presenting >12-24 hours after injury (which should heal by secondary intention).

\section*{Risks}
Wound infection, dehiscence (opening of the wound), scarring, or local nerve injury.

\section*{Aftercare and Recovery}
Apply antibiotic ointment and a sterile dressing. Keep the wound clean and dry for 24-48 hours.

\section*{Results and Interpretation}
Successful closure shows aligned skin edges, minimal tension, and absence of active bleeding or infection.

\section*{Expected Results}
Well-approximated wound edges; no active bleeding; minimal swelling.

\section*{Follow-up}
Suture removal in 5-7 days (face), 7-10 days (scalp/trunk), or 10-14 days (extremities).

\section*{References}
\begin{enumerate}
  \item MedlinePlus. Wound Suturing. U.S. National Library of Medicine; 2025.
  \item Current Medical Diagnosis and Treatment. McGraw-Hill; 2025.
\end{enumerate}

\clearpage
